% Figure: the third runtime - state written by the RUNTIME, not by the model.
%
% Same visual language as fig1_architecture (box column, group frame, growth
% panel) so the three designs can be compared at a glance. The point of the
% figure is WHERE the state block comes from: in Figure 1 the model emits it as
% generated text, here the runtime derives it from what already happened.
%
% Contains no measured numbers; the measured ladder is the composition figure.

\begin{figure*}[t]
\centering
\resizebox{\textwidth}{!}{%
\begin{tikzpicture}[
  font=\small,
  box/.style={draw, rounded corners=2pt, minimum width=3.4cm,
              minimum height=0.6cm, align=center, inner sep=2pt},
  hist/.style={box, fill=red!8},
  state/.style={box, fill=blue!10, draw=blue!60, thick},
  rtstate/.style={box, fill=green!12, draw=green!55!black, thick},
  plain/.style={box, fill=black!3},
  llm/.style={box, fill=black!8, minimum height=0.85cm, thick},
  rt/.style={box, fill=green!8, draw=green!55!black, minimum height=0.85cm, thick},
  grp/.style={draw, dashed, rounded corners=3pt, inner sep=7pt},
  arr/.style={-latex, thick},
]

\node[font=\bfseries] at (0,4.55) {Состояние, написанное моделью (\SkillState{})};
\node[plain]  (p1) at (0, 3.05) {спецификация задачи $P$};
\node[state]  (s1) at (0, 2.35) {состояние выполнения $\Sigma_{t-1}$};
\node[plain]  (o1) at (0, 1.65) {последнее наблюдение $o_t$};
\node[grp, fit=(p1)(s1)(o1)] (g1) {};
\node[above=2pt of g1, font=\bfseries\footnotesize, text=blue!60!black]
  {промпт ограничен};
\node[llm] (m1) at (0,0.35) {языковая модель};
\draw[arr] (g1.south) -- (m1.north);
\node[plain, dashed] (a1) at (0,-0.85) {выполнить действие $a_t$};
\draw[arr] (m1) -- (a1);
\draw[arr, blue!60!black] (m1.west) -- ++(-1.3,0) |- (s1.west)
  node[pos=0.26, left, font=\scriptsize, text=blue!60!black, align=center]
  {модель выдаёт\\$\Delta\Sigma_t$ как текст};
\node[font=\scriptsize, text=red!65!black, align=center] at (0,-1.75)
  {состояние стоит выходных токенов;\\модель может перезаписать свои факты};

\begin{scope}[xshift=7.8cm]
\node[font=\bfseries] at (0,4.55) {Состояние, написанное рантаймом (эта работа)};
\node[plain]   (p2) at (0, 3.05) {спецификация задачи $P$};
\node[rtstate] (s2) at (0, 2.35) {проверенное состояние $V_{t-1}$};
\node[plain]   (o2) at (0, 1.65) {последнее наблюдение $o_t$};
\node[grp, fit=(p2)(s2)(o2)] (g2) {};
\node[above=2pt of g2, font=\bfseries\footnotesize, text=green!45!black]
  {промпт ограничен};
\node[llm] (m2) at (0,0.35) {языковая модель};
\draw[arr] (g2.south) -- (m2.north);
\node[plain, dashed] (a2) at (0,-0.85) {выполнить действие $a_t$};
\draw[arr] (m2) -- (a2);

\node[rt] (rec) at (4.6,-0.85) {регистратор рантайма};
\draw[arr, green!50!black] (a2.east) -- (rec.west)
  node[midway, below=1pt, font=\scriptsize, text=green!45!black, align=center]
  {результат\\инструмента};
\draw[arr, green!50!black] (rec.north) -- ++(0,1.45) -| (s2.east)
  node[pos=0.28, right, font=\scriptsize, text=green!45!black, align=left]
  {выведено, а не\\сгенерировано};
\node[font=\scriptsize, text=green!35!black, align=center] at (0,-1.75)
  {состояние не стоит выходных токенов;\\модель не может его перезаписать};
\end{scope}

\draw[dashed, black!30] (-3.4,-2.45) -- (11.2,-2.45);

\node[font=\bfseries\small, align=center] at (3.9,-3.05)
  {Что рантайм может записать, не спрашивая модель};
\begin{scope}[yshift=-3.95cm]
  \foreach \x/\tag/\what in {
     -3.0/v1/{затронутые файлы},
     -0.6/v2/{+ содержимое},
      1.8/v3/{+ команды},
      4.2/v4/{+ диагностика},
      6.6/v5/{+ окружение},
      9.0/v6/{+ тесты},
     11.4/v7/{+ провалы}} {
    \node[draw, rounded corners=2pt, fill=green!10, draw=green!55!black,
          minimum width=2.05cm, minimum height=0.62cm, font=\scriptsize,
          align=center] at (\x,0) {\textbf{\tag}\\\what};
  }
  \draw[arr, green!50!black] (-4.2,-0.68) -- (12.5,-0.68)
    node[midway, below, font=\scriptsize, text=green!45!black]
    {больше делегировано рантайму, один компонент на ступень};
\end{scope}

\end{tikzpicture}}
\caption{Два способа ограничить промпт блоком состояния, различающихся тем,
кто его пишет. \textbf{Слева:} модель выдаёт обновление состояния как
сгенерированный текст, так что его поддержание тарифицируется по выходному
тарифу, а модель вольна перезаписать ранее записанный факт, чему оригинал
приписывает \SI{68}{\percent} своих отказов моделей с открытыми
весами\origin. \textbf{Справа:} рантайм выводит ту же запись из результатов
инструментов, которые уже произошли. Блок — это факт о том, что случилось, а
не утверждение модели об этом, он не стоит сгенерированных токенов и не
может быть перезаписан моделью. Полоса внизу — это фактор, варьируемый в
\S\ref{sec:composition}: каждая ступень передаёт рантайму ещё один
компонент, при замороженных хвосте и схеме, так что разрыв между соседями
называет один делегированный компонент.}
\label{fig:runtime-written}
\end{figure*}
